%=============================================================
%  LIBRO: Criptografia de Reticulos y Protocolo de Votacion
%         Electronica Post-Cuantica
%  Proyecto de Grado — UMSA, 2026
%=============================================================
\documentclass[12pt,a4paper,twoside,openright]{book}

%------------------------------------------------------------
% PAQUETES BASICOS
%------------------------------------------------------------
\usepackage[utf8]{inputenc}
\usepackage[T1]{fontenc}
\usepackage[spanish,es-noquoting]{babel}
\usepackage{lmodern}
\usepackage{microtype}

%------------------------------------------------------------
% MATEMATICAS
%------------------------------------------------------------
\usepackage{amsmath,amssymb,amsthm,mathtools}
\usepackage{bm}

%------------------------------------------------------------
% GEOMETRIA
%------------------------------------------------------------
\usepackage[
  a4paper,
  top=2.5cm, bottom=2.5cm,
  inner=3cm, outer=2cm,
  headheight=15pt
]{geometry}

%------------------------------------------------------------
% CABECERAS
%------------------------------------------------------------
\usepackage{fancyhdr}
\pagestyle{fancy}
\fancyhf{}
\fancyhead[LE]{\small\leftmark}
\fancyhead[RO]{\small\rightmark}
\fancyfoot[C]{\thepage}
\renewcommand{\headrulewidth}{0.4pt}

%------------------------------------------------------------
% COLORES Y CAJAS
%------------------------------------------------------------
\usepackage[dvipsnames,svgnames,x11names]{xcolor}
\usepackage{tcolorbox}
\tcbuselibrary{breakable,skins}

\definecolor{PQCBlue}{RGB}{25,60,120}
\definecolor{PQCCyan}{RGB}{0,150,180}
\definecolor{PQCGreen}{RGB}{30,130,80}
\definecolor{PQCOrange}{RGB}{200,100,20}
\definecolor{PQCGray}{RGB}{95,95,95}
\definecolor{BoxBg}{RGB}{245,248,255}
\definecolor{WarnBg}{RGB}{255,248,235}

\newtcolorbox{midef}[2][]{
  colback=BoxBg, colframe=PQCBlue, fonttitle=\bfseries,
  title={Definici\'{o}n: #2}, #1, breakable
}
\newtcolorbox{miteo}[2][]{
  colback=BoxBg, colframe=PQCGreen, fonttitle=\bfseries,
  title={Teorema: #2}, #1, breakable
}
\newtcolorbox{minota}[1][Nota]{
  colback=WarnBg, colframe=PQCOrange, fonttitle=\bfseries,
  title={#1}, breakable
}

%------------------------------------------------------------
% CODIGO
%------------------------------------------------------------
\usepackage{listings}
\lstdefinestyle{PQCstyle}{
  backgroundcolor=\color{gray!8},
  basicstyle=\ttfamily\footnotesize,
  keywordstyle=\color{PQCBlue}\bfseries,
  commentstyle=\color{PQCGray}\itshape,
  stringstyle=\color{PQCGreen},
  numbers=left,
  numberstyle=\tiny\color{PQCGray},
  breaklines=true,
  frame=single,
  framerule=0.4pt,
  rulecolor=\color{PQCBlue!40},
  captionpos=b,
  tabsize=4,
  showstringspaces=false,
}
\lstset{style=PQCstyle}

%------------------------------------------------------------
% TABLAS Y FIGURAS
%------------------------------------------------------------
\usepackage{booktabs}
\usepackage{multirow}
\usepackage{array}
\usepackage{longtable}
\usepackage{graphicx}
\usepackage{float}
\usepackage{caption}
\captionsetup{font=small, labelfont=bf}

%------------------------------------------------------------
% ALGORITMOS
%------------------------------------------------------------
\usepackage[ruled,vlined,linesnumbered]{algorithm2e}
\SetAlFnt{\small}
\SetAlCapFnt{\small}

%------------------------------------------------------------
% ENLACES
%------------------------------------------------------------
\usepackage[
  colorlinks=true,
  linkcolor=PQCBlue,
  urlcolor=PQCCyan,
  citecolor=PQCGreen,
  pdftitle={Criptografia de Reticulos y Protocolo de Votacion PQC},
  pdfauthor={Proyecto de Grado UMSA 2026},
  pdfsubject={Criptografia Post-Cuantica}
]{hyperref}

%------------------------------------------------------------
% INDICE
%------------------------------------------------------------
\usepackage{makeidx}
\makeindex

%------------------------------------------------------------
% ENTORNOS MATEMATICOS
%------------------------------------------------------------
\theoremstyle{plain}
\newtheorem{thm}{Teorema}[chapter]
\newtheorem{lem}[thm]{Lema}
\newtheorem{prop}[thm]{Proposici\'{o}n}
\newtheorem{cor}[thm]{Corolario}

\theoremstyle{definition}
\newtheorem{defn}[thm]{Definici\'{o}n}
\newtheorem{ej}[thm]{Ejemplo}

\theoremstyle{remark}
\newtheorem{obs}[thm]{Observaci\'{o}n}

%------------------------------------------------------------
% COMANDOS MATEMATICOS
%------------------------------------------------------------
\newcommand{\Z}{\mathbb{Z}}
\newcommand{\Zq}{\mathbb{Z}_q}
\newcommand{\R}{\mathbb{R}}
\newcommand{\Rq}{R_q}
\newcommand{\vv}[1]{\mathbf{#1}}
\newcommand{\mm}[1]{\mathbf{#1}}
\newcommand{\norm}[1]{\left\|#1\right\|}
\newcommand{\inner}[2]{\langle #1,\, #2 \rangle}
\newcommand{\MLKEM}{\textsf{ML-KEM}}
\newcommand{\MLDSA}{\textsf{ML-DSA}}
\newcommand{\KeyGen}{\mathsf{KeyGen}}
\newcommand{\Encaps}{\mathsf{Encaps}}
\newcommand{\Decaps}{\mathsf{Decaps}}
\newcommand{\Sign}{\mathsf{Sign}}
\newcommand{\Verify}{\mathsf{Verify}}
\newcommand{\negl}{\mathsf{negl}}

%=============================================================
\begin{document}
%=============================================================

%------------------------------------------------------------
% PORTADA
%------------------------------------------------------------
\begin{titlepage}
\centering
\vspace*{1cm}
{\huge\bfseries\color{PQCBlue}
  Criptograf\'ia de Ret\'iculos\\[0.4em]
  y Protocolo de Votaci\'on\\[0.4em]
  Electr\'onica Post-Cu\'antica
\par}
\vspace{0.5cm}
\rule{\linewidth}{0.5mm}
\vspace{0.8cm}

{\Large\itshape
  Una introducci\'on t\'ecnica a los fundamentos matem\'aticos\\
  de los ret\'iculos, los algoritmos \MLKEM{} y \MLDSA{},\\
  y su aplicaci\'on en un sistema de votaci\'on electr\'onica
\par}
\vspace{1.2cm}

{\large\bfseries Proyecto de Grado}\\[0.3em]
{\large Universidad Mayor de San Andr\'es (UMSA)}\\[0.2em]
{\large Facultad de Ciencias Puras y Naturales}\\[0.2em]
{\large Carrera de Inform\'atica}\\[0.8em]

\rule{0.6\linewidth}{0.3mm}\\[0.5em]
{\large\textbf{Metodolog\'ia:} Design Science Research (DSR)}\\[0.2em]
{\large\textbf{Est\'andares:} NIST FIPS 203 / FIPS 204 (2024)}\\[0.2em]
{\large\textbf{A\~no:} 2026}

\vfill
{\small\color{PQCGray}
  \textit{``La seguridad futura depende de fundamentos matem\'aticos\\
  que ni las computadoras cu\'anticas puedan quebrantar.''}
}
\end{titlepage}

%------------------------------------------------------------
\frontmatter
%------------------------------------------------------------
\tableofcontents
\listoftables

%------------------------------------------------------------
% PREFACIO
%------------------------------------------------------------
\chapter*{Prefacio}
\addcontentsline{toc}{chapter}{Prefacio}

Este libro nace como compa\~nero t\'ecnico del proyecto de grado
\textit{``Protocolo para la votaci\'on electr\'onica basado en criptograf\'ia
post-cu\'antica de ret\'iculos''}, desarrollado en la Universidad Mayor de San
Andr\'es (UMSA) durante el a\~no 2026 bajo la metodolog\'ia \textit{Design
Science Research} (DSR).

\medskip
La motivaci\'on principal es doble. Por un lado, los avances en computaci\'on
cu\'antica ---en particular el algoritmo de Shor (1994)--- amenazan romper los
criptosistemas cl\'asicos basados en la factorizaci\'on de enteros (RSA) y el
problema del logaritmo discreto (ECDSA). Por el otro, los sistemas de
votaci\'on electr\'onica son infraestructura cr\'itica cuya integridad debe
garantizarse no s\'olo hoy, sino durante d\'ecadas.

\medskip
El libro est\'a pensado para un lector con conocimientos b\'asicos de \'algebra
lineal y aritm\'etica modular, pero que \emph{se est\'a iniciando} en la
criptograf\'ia de ret\'iculos y en los est\'andares post-cu\'anticos del NIST.

\medskip
La estructura sigue una progresi\'on natural:
\begin{enumerate}
  \item \textbf{Contexto hist\'orico} --- por qu\'e la criptograf\'ia cl\'asica
        est\'a en riesgo.
  \item \textbf{Matem\'aticas de ret\'iculos} --- definiciones, problemas
        dif\'iciles y construcciones algebraicas.
  \item \textbf{Primitivas criptogr\'aficas} --- LWE, MLWE, esquemas KEM y DSA.
  \item \textbf{El protocolo de votaci\'on} --- dise\~no, implementaci\'on y
        an\'alisis emp\'irico.
  \item \textbf{Resultados} --- benchmarks y conclusiones.
\end{enumerate}

\medskip
La gu\'ia de Tamarin Prover existe como documento separado
(\texttt{tamarin/guia\_tamarin\_libro.tex}); este libro ofrece s\'olo
una referencia m\'inima a los lemas verificados.

\vspace{1cm}
\noindent\textit{UMSA, La Paz --- 2026}

%=============================================================
\mainmatter
%=============================================================

%###############################################################
\part{Fundamentos Matem\'aticos y Criptogr\'aficos}
%###############################################################

%==============================================================
\chapter{El Desaf\'io Cu\'antico a la Criptograf\'ia Cl\'asica}
%==============================================================
\label{chap:desafio}

\section{Historia breve de la criptograf\'ia de clave p\'ublica}

La criptograf\'ia de clave p\'ublica naci\'o en 1976 con el art\'iculo
seminal de Diffie y Hellman, \textit{New Directions in Cryptography}. La idea
revolucionaria era que dos partes pod\'ian establecer un secreto compartido
a trav\'es de un canal inseguro sin haberse reunido previamente.

\medskip
Durante los a\~nos siguientes surgieron tres familias dominantes:

\begin{description}
  \item[RSA (1978)] La seguridad descansa en la \emph{dificultad de factorizar}
    el producto de dos primos grandes $p$ y $q$. La clave p\'ublica es $n = pq$
    y la privada es el par $(p, q)$.
  \item[Diffie-Hellman / ElGamal] La seguridad se basa en el
    \emph{problema del logaritmo discreto} en grupos finitos $\Z_p^*$ o
    curvas el\'ipticas.
  \item[ECDSA / EdDSA] Firmas digitales sobre curvas el\'ipticas. ECDSA P-256
    ofrece $\sim$128 bits de seguridad con claves muy compactas ($\sim$64 bytes).
\end{description}

\subsection{Por qu\'e estos sistemas est\'an en riesgo}

En 1994, Peter Shor public\'o un algoritmo cu\'antico de tiempo polinomial
que resuelve tanto la factorizaci\'on de enteros como el logaritmo discreto.
Sobre una computadora cu\'antica universal con suficientes qubits l\'ogicos:

\begin{itemize}
  \item Factoriza $n = pq$ en $O((\log n)^3)$ operaciones cu\'anticas.
  \item Resuelve el logaritmo discreto en grupos finitos y en curvas el\'ipticas.
\end{itemize}

\begin{minota}[\textquestiondown Cu\'ando ocurrir\'a?]
  A la fecha de este proyecto (2026) no existe un computador cu\'antico
  universal con los miles de qubits l\'ogicos necesarios para romper RSA-2048.
  Sin embargo, el principio de \emph{``harvest now, decrypt later''} implica
  que los adversarios pueden recopilar tr\'afico cifrado hoy y descifrarlo
  en el futuro. Por ello, la migraci\'on a PQC es urgente.
\end{minota}

\section{La respuesta: Criptograf\'ia Post-Cu\'antica}

La \textbf{Criptograf\'ia Post-Cu\'antica (PQC)} es el conjunto de algoritmos
cuya seguridad se basa en problemas dif\'iciles incluso para computadoras
cu\'anticas. En 2016, el NIST inici\'o un proceso de estandarizaci\'on que
culmin\'o en 2024:

\begin{table}[H]
\centering
\caption{Est\'andares PQC del NIST (2024)}
\label{tab:nist_standards}
\begin{tabular}{llll}
\toprule
\textbf{FIPS} & \textbf{Algoritmo} & \textbf{Funci\'on} & \textbf{Base matem\'atica} \\
\midrule
203 & ML-KEM & Encapsulamiento de clave & Module-LWE \\
204 & ML-DSA & Firma digital & Module-LWE (Dilithium) \\
205 & SLH-DSA & Firma digital & Hash-based (SPHINCS+) \\
\bottomrule
\end{tabular}
\end{table}

Este proyecto implementa \textbf{ML-KEM-768} (FIPS 203) y
\textbf{ML-DSA-65} (FIPS 204), ambos basados en Module-LWE.

%==============================================================
\chapter{Ret\'iculos: Definiciones y Geometr\'ia}
%==============================================================
\label{chap:reticulos}

\section{\textquestiondown Qu\'e es un ret\'iculo?}

\begin{midef}{Ret\'iculo (Lattice)}
  Sea $\vv{b}_1, \ldots, \vv{b}_n \in \R^m$ un conjunto de vectores
  linealmente independientes ($n \leq m$). El \emph{ret\'iculo}
  $\mathcal{L}$ generado por esta base es:
  \[
    \mathcal{L}(\vv{b}_1,\ldots,\vv{b}_n)
    = \left\{ \sum_{i=1}^{n} z_i \, \vv{b}_i \;\middle|\; z_i \in \Z \right\}.
  \]
  La matriz $\mm{B} = [\vv{b}_1 | \cdots | \vv{b}_n]$ se llama
  \emph{base} del ret\'iculo.
\end{midef}

Un ret\'iculo es, intuitivamente, una \emph{cuadr\'icula generalizada} en el
espacio $n$-dimensional: un conjunto discreto e infinito de puntos con
estructura peri\'odica.

\subsection{Ejemplos b\'asicos}

\begin{ej}[Ret\'iculo entero $\Z^2$]
  Tomando $\vv{b}_1 = (1,0)$ y $\vv{b}_2 = (0,1)$, obtenemos
  $\mathcal{L} = \Z^2$: la cuadr\'icula entera del plano.
\end{ej}

\begin{ej}[Ambig\"uedad de la base]
  $\vv{b}_1' = (2,1)$ y $\vv{b}_2' = (0,1)$ generan el mismo ret\'iculo
  $\Z^2$. Esta ambig\"uedad es clave en criptograf\'ia: la \emph{clave
  p\'ublica} es una base ``mala'' (vectores largos) y la \emph{clave
  privada} es una base ``buena'' (vectores cortos, f\'acil de usar para
  decodificar).
\end{ej}

\section{Propiedades geom\'etricas fundamentales}

\subsection{Determinante y volumen fundamental}

\begin{defn}
  El \emph{determinante} de un ret\'iculo con base $\mm{B} \in \R^{n\times n}$
  es $\det(\mathcal{L}) = |\det(\mm{B})|$. Geom\'etricamente, es el volumen
  del dominio fundamental (celda unitaria).
\end{defn}

\subsection{Primer m\'inimo}

\begin{defn}
  El \emph{primer m\'inimo} $\lambda_1(\mathcal{L})$ es la norma del vector
  m\'as corto no nulo:
  \[
    \lambda_1(\mathcal{L}) = \min_{\vv{v} \in \mathcal{L},\, \vv{v}\neq\vv{0}}
    \|\vv{v}\|.
  \]
\end{defn}

La cota de Hermite establece:
$\lambda_1(\mathcal{L}) \leq \sqrt{n} \cdot \det(\mathcal{L})^{1/n}$.

\section{Problemas dif\'iciles en ret\'iculos}

Los criptosistemas de ret\'iculos se basan en problemas sin algoritmo
eficiente cl\'asico \emph{ni cu\'antico} conocido.

\subsection{SVP --- Problema del Vector M\'as Corto}

\begin{midef}{SVP (Shortest Vector Problem)}
  Dado un ret\'iculo $\mathcal{L}$ por una base arbitraria, encontrar
  $\vv{v} \in \mathcal{L}$, $\vv{v}\neq\vv{0}$, tal que
  $\|\vv{v}\| = \lambda_1(\mathcal{L})$.
\end{midef}

SVP exacto es NP-dif\'icil bajo reducciones aleatorias. En la pr\'actica se
trabaja con la versi\'on aproximada: dado un ret\'iculo, decidir si
$\lambda_1 \leq 1$ o $\lambda_1 > \gamma$ para un factor de aproximaci\'on
$\gamma$.

\subsection{CVP --- Problema del Vector M\'as Cercano}

\begin{midef}{CVP (Closest Vector Problem)}
  Dado un ret\'iculo $\mathcal{L}$ y un punto objetivo $\vv{t} \in \R^m$,
  encontrar $\vv{v}^* = \arg\min_{\vv{v}\in\mathcal{L}} \|\vv{v} - \vv{t}\|$.
\end{midef}

CVP es NP-dif\'icil y es la base de los mecanismos de \emph{decodificaci\'on}
en los esquemas de encapsulamiento.

\subsection{SIS --- Short Integer Solution}

\begin{midef}{SIS (Short Integer Solution)}
  Dados vectores uniformes $\vv{a}_1,\ldots,\vv{a}_m \in \Zq^n$
  (en la matriz $\mm{A} \in \Zq^{n\times m}$), encontrar
  $\vv{z} \in \Z^m$, $\vv{z}\neq\vv{0}$, con $\|\vv{z}\| \leq \beta$ y
  $\mm{A}\vv{z} = \vv{0} \pmod{q}$.
\end{midef}

SIS es la base de los esquemas de \emph{firma digital} basados en ret\'iculos
(incluyendo ML-DSA).

\begin{table}[H]
\centering
\caption{Complejidad de los mejores algoritmos para SVP en rango $n$}
\label{tab:complejidad_svp}
\begin{tabular}{lll}
\toprule
\textbf{Algoritmo} & \textbf{Complejidad (tiempo)} & \textbf{Tipo} \\
\midrule
LLL (1982) & $\mathrm{poly}(n)$ & Cl\'asico (aproximado $2^{O(n)}$) \\
BKZ (1987) & $2^{O(n/\log n)}$ & Cl\'asico \\
Sieving (Ajtai+) & $2^{O(n)}$ & Cl\'asico/Cu\'antico \\
Algoritmo de Shor & N/A & No aplica a SVP \\
\bottomrule
\end{tabular}
\end{table}

Ning\'un algoritmo cu\'antico conocido proporciona una aceleraci\'on
superpolinomial sobre los mejores algoritmos cl\'asicos para SVP/CVP, lo
que justifica la seguridad post-cu\'antica de los criptosistemas de ret\'iculos.

%==============================================================
\chapter{Learning With Errors (LWE) y sus Variantes}
%==============================================================
\label{chap:lwe}

\section{El Problema LWE}

Introducido por Regev en 2005, \textbf{LWE} (Learning With Errors) es la
piedra angular de la criptograf\'ia de ret\'iculos moderna. La idea es que
esconder un secreto en un sistema lineal \emph{con errores peque\~nos}
resulta computacionalmente intratable.

\begin{midef}{LWE --- Distribuci\'on}
  Sea $n$ la dimensi\'on de seguridad, $q$ el m\'odulo (primo), y
  $\chi$ una distribuci\'on de error (Gaussiana discreta centrada). La
  \emph{distribuci\'on LWE} $A_{\vv{s},\chi}$ para secreto
  $\vv{s} \in \Zq^n$ genera muestras $(\vv{a}, b) \in \Zq^n \times \Zq$:
  \[
    \vv{a} \xleftarrow{\$} \Zq^n, \quad
    e \leftarrow \chi, \quad
    b = \inner{\vv{a}}{\vv{s}} + e \pmod{q}.
  \]
\end{midef}

\begin{midef}{Problema de decisi\'on LWE}
  Dadas $m$ muestras, distinguir si provienen de $A_{\vv{s},\chi}$
  (LWE con secreto $\vv{s}$) o de la distribuci\'on uniforme
  $\mathcal{U}(\Zq^n \times \Zq)$, con ventaja no despreciable.
\end{midef}

\begin{miteo}{Dificultad cu\'antica de LWE}
  Bajo reducciones cu\'anticas, resolver LWE (incluso en promedio) es al
  menos tan dif\'icil como resolver GapSVP$_\gamma$ en el \emph{peor caso}
  (Regev 2005, Peikert 2009).
\end{miteo}

\subsection{Intuici\'on: por qu\'e LWE es dif\'icil}

Sin el error $e$, el sistema $b_i = \inner{\vv{a}_i}{\vv{s}} \pmod{q}$ es
resoluble en tiempo polinomial con eliminaci\'on gaussiana. El error convierte
el sistema en uno \emph{sobredeterminado con ruido}: los algoritmos
especializados (BKZ, sieving) requieren tiempo exponencial.

\section{Ring-LWE (RLWE)}

LWE est\'andar requiere almacenar matrices de tama\~no $O(n^2)$, lo que
resulta ineficiente. \textbf{RLWE} (Lyubashevsky, Peikert, Regev 2010)
traslada el problema al anillo de polinomios c\'iclicos:

\begin{midef}{Anillo c\'iclico $R_q$}
  \[
    \Rq = \Z_q[x] / (x^n + 1)
  \]
  donde $n$ es una potencia de 2. Los elementos son polinomios de grado
  $< n$ con coeficientes en $\Zq$.
\end{midef}

\subsection{Ventajas de $R_q$}

\begin{enumerate}
  \item La multiplicaci\'on polinomial en $\Rq$ se puede computar en
    $O(n \log n)$ usando la \textbf{NTT} (Number Theoretic Transform),
    an\'aloga a la DFT pero en aritm\'etica modular.
  \item La clave p\'ublica es un \'unico polinomio en lugar de una
    matriz completa: reducci\'on de $O(n^2)$ a $O(n)$ en tama\~no.
\end{enumerate}

\section{Module-LWE (MLWE) --- La base de ML-KEM y ML-DSA}

\textbf{MLWE} (Langlois y Stehl\'e, 2012) generaliza LWE y RLWE usando
\emph{m\'odulos} de anillos, equilibrando eficiencia y confianza en la
seguridad.

\begin{midef}{Module-LWE}
  Sea $k \geq 1$ el rango del m\'odulo. En MLWE se trabaja con
  vectores $\vv{s} \in R_q^k$ (secreto) y matrices $\mm{A} \in R_q^{k \times k}$
  (p\'ublicas). Una instancia MLWE genera:
  \[
    \mm{A} \xleftarrow{\$} R_q^{k\times k}, \quad
    \vv{s}, \vv{e} \leftarrow \chi^k, \quad
    \vv{b} = \mm{A}\vv{s} + \vv{e} \in R_q^k.
  \]
  El problema de decisi\'on MLWE es distinguir $(\mm{A}, \vv{b})$ de
  $(\mm{A}, \vv{u})$ con $\vv{u} \xleftarrow{\$} R_q^k$.
\end{midef}

\subsection{Par\'ametros de ML-KEM-768}

\begin{table}[H]
\centering
\caption{Par\'ametros criptogr\'aficos de ML-KEM-768 (FIPS 203)}
\label{tab:params_mlkem768}
\begin{tabular}{lll}
\toprule
\textbf{Par\'ametro} & \textbf{Valor} & \textbf{Descripci\'on} \\
\midrule
$n$ & 256 & Grado del polinomio c\'iclico \\
$q$ & 3329 & M\'odulo primo \\
$k$ & 3 & Rango del m\'odulo \\
$\eta_1$ & 2 & Par\'ametro de error para secreto \\
$\eta_2$ & 2 & Par\'ametro de error para encapsulamiento \\
$(d_u, d_v)$ & (10, 4) & Bits de compresi\'on del ciphertext \\
\midrule
\multicolumn{3}{l}{\textbf{Tama\~nos}} \\
Clave p\'ublica ($\mathtt{ek}$) & 1\,184 bytes & Incluye semilla de $\mm{A}$ y $\vv{t}$ \\
Clave privada ($\mathtt{dk}$) & 2\,400 bytes & Incluye $\vv{s}$, hash de $\mathtt{ek}$, $z$ \\
Ciphertext ($\mathtt{ct}$) & 1\,088 bytes & Encapsulamiento de $K$ \\
\bottomrule
\end{tabular}
\end{table}

\subsection{La Transformada NTT}

Con $q \equiv 1 \pmod{2n}$, existen ra\'ices $n$-\'esimas de la unidad en
$\Zq$. La NTT permite multiplicar polinomios eficientemente:
\begin{enumerate}
  \item Transformar $a(x), b(x)$ al dominio NTT: $\hat{a}, \hat{b} \in \Zq^n$.
  \item Multiplicar puntualmente: $\hat{c}_i = \hat{a}_i \cdot \hat{b}_i$.
  \item Aplicar NTT inversa: $c(x) = \mathrm{INTT}(\hat{c})$.
\end{enumerate}
Costo total: $O(n \log n)$ en lugar de $O(n^2)$ con convo-luci\'on directa.

%==============================================================
\chapter{Esquemas de Encapsulamiento de Clave (KEM)}
%==============================================================
\label{chap:kem}

\section{Concepto de KEM}

Un \textbf{KEM} (Key Encapsulation Mechanism) es un protocolo asim\'etrico
que permite a dos partes establecer un secreto compartido $K$ sin
transmitirlo directamente. Tiene tres algoritmos:

\begin{description}
  \item[$\KeyGen() \to (\mathtt{ek}, \mathtt{dk})$] Genera un par de claves:
    p\'ublica de encapsulamiento (\texttt{ek}) y privada de
    desencapsulamiento (\texttt{dk}).
  \item[$\Encaps(\mathtt{ek}) \to (\mathtt{ct}, K)$] Usando la clave p\'ublica,
    genera un texto cifrado $\mathtt{ct}$ y un secreto compartido $K$.
  \item[$\Decaps(\mathtt{dk}, \mathtt{ct}) \to K$] Usando la clave privada,
    recupera el secreto $K$ a partir de $\mathtt{ct}$.
\end{description}

\begin{minota}
  En un KEM seguro: (1) \textbf{Correcci\'on}: $\Decaps(\mathtt{dk},
  \Encaps(\mathtt{ek})) = K$; (2) \textbf{Seguridad IND-CCA2}: $K$ es
  indistinguible de un valor aleatorio para cualquier adversario que no
  posea $\mathtt{dk}$.
\end{minota}

\section{Construcci\'on de ML-KEM (FIPS 203)}

ML-KEM es el KEM estandarizado por el NIST basado en MLWE. Internamente es
una versi\'on CCA-segura del esquema Kyber, transformada mediante la t\'ecnica
de Fujisaki-Okamoto (FO).

\subsection{Generaci\'on de claves}

\begin{algorithm}[H]
\caption{$\MLKEM.\KeyGen()$}
\label{alg:mlkem_keygen}
\SetKwInOut{Output}{Salida}
\Output{Par de claves $(\mathtt{ek}, \mathtt{dk})$}
Muestrear semilla aleatoria $d \xleftarrow{\$} \{0,1\}^{32}$\;
$(d', z) \leftarrow H(d)$\tcp*{H: hash extendible (XOF)}
$(\hat{\mm{A}}, \vv{s}, \vv{e}) \leftarrow \mathrm{SampleNoise}(d')$\;
$\vv{t} \leftarrow \hat{\mm{A}}\hat{\vv{s}} + \hat{\vv{e}}$\tcp*{en dominio NTT}
$\mathtt{ek} \leftarrow (\vv{t},\; \mathrm{seed}(\mm{A}))$\tcp*{clave p\'ublica}
$\mathtt{dk} \leftarrow (\hat{\vv{s}},\; \mathtt{ek},\; H(\mathtt{ek}),\; z)$\tcp*{clave privada}
\Return $(\mathtt{ek}, \mathtt{dk})$
\end{algorithm}

\subsection{Encapsulamiento}

\begin{algorithm}[H]
\caption{$\MLKEM.\Encaps(\mathtt{ek})$}
\label{alg:mlkem_encaps}
\SetKwInOut{Output}{Salida}
\Output{$(K, \mathtt{ct})$}
$m \xleftarrow{\$} \{0,1\}^{32}$\tcp*{valor aleatorio ef\'imero}
$(K, r) \leftarrow G(m,\; H(\mathtt{ek}))$\tcp*{G: funci\'on hash}
$(\vv{u}, v) \leftarrow \mathrm{PKE.Enc}(\mathtt{ek},\; m;\; r)$\tcp*{cifrado PKE subyacente}
$\mathtt{ct} \leftarrow \mathrm{Compress}(\vv{u}, v)$\;
\Return $(K, \mathtt{ct})$
\end{algorithm}

El PKE subyacente opera sobre MLWE:
\begin{align}
  \vv{r} &\leftarrow \chi^k, \quad \vv{e}_1 \leftarrow \chi^k, \quad e_2 \leftarrow \chi \notag\\
  \vv{u} &= \mm{A}^T \vv{r} + \vv{e}_1 \pmod{q}\\
  v      &= \vv{t}^T \vv{r} + e_2 + \lfloor q/2 \rfloor \cdot m \pmod{q}
\end{align}

\subsection{Desencapsulamiento}

\begin{algorithm}[H]
\caption{$\MLKEM.\Decaps(\mathtt{dk}, \mathtt{ct})$}
\label{alg:mlkem_decaps}
\SetKwInOut{Output}{Salida}
\Output{Secreto compartido $K$}
$(\vv{u}', v') \leftarrow \mathrm{Decompress}(\mathtt{ct})$\;
$m' \leftarrow \mathrm{Round}(v' - \hat{\vv{s}}^T \vv{u}')$\tcp*{recuperar $m$}
$(K', r') \leftarrow G(m',\; H(\mathtt{ek}))$\;
Reencapsular: $\mathtt{ct}' \leftarrow \mathrm{PKE.Enc}(\mathtt{ek},\; m';\; r')$\;
\If{$\mathtt{ct}' = \mathtt{ct}$}{
  \Return $K'$\;
}
\Else{
  \Return $\mathrm{PRF}(z,\; \mathtt{ct})$\tcp*{rechazo impl\'icito}
}
\end{algorithm}

\subsection{Por qu\'e el desencapsulamiento funciona}

La operaci\'on clave es:
\[
  v' - \vv{s}^T \vv{u}'
  = \vv{t}^T\vv{r} + e_2 + \tfrac{q}{2}m - \vv{s}^T(\mm{A}^T\vv{r} + \vv{e}_1)
  = \tfrac{q}{2}m + \underbrace{e_2 + \vv{e}^T\vv{r} - \vv{s}^T\vv{e}_1}_{\text{ruido peque\~no}}
\]
Si el ruido total es menor que $q/4$, el redondeo a $\{0,1\}$ recupera $m$
correctamente. La verificaci\'on de re-encapsulamiento garantiza IND-CCA2
(transformaci\'on Fujisaki-Okamoto).

%==============================================================
\chapter{Firmas Digitales Basadas en Ret\'iculos}
%==============================================================
\label{chap:firmas}

\section{Concepto de firma digital}

Una firma digital permite verificar que un mensaje $m$ fue producido por
quien posee la clave privada, y que no ha sido modificado:

\begin{description}
  \item[$\KeyGen() \to (\mathtt{sk}, \mathtt{pk})$] Genera el par de claves.
  \item[$\Sign(\mathtt{sk}, m) \to \sigma$] Produce una firma $\sigma$.
  \item[$\Verify(\mathtt{pk}, m, \sigma) \to \{0,1\}$] Verifica la firma.
\end{description}

La propiedad de seguridad central es la \textbf{EUF-CMA}: un adversario que
puede solicitar firmas de mensajes de su elecci\'on no puede forjar una firma
v\'alida sobre ning\'un mensaje nuevo.

\section{Construcci\'on de ML-DSA (FIPS 204) --- Dilithium}

ML-DSA se basa en el esquema Dilithium (Ducas et al., 2018), que usa el
paradigma \emph{Fiat-Shamir con abortos}.

\subsection{Idea intuitiva}

El protocolo de identificaci\'on subyacente funciona as\'i:
\begin{enumerate}
  \item El firmante elige $\vv{y} \xleftarrow{\$} R_q^k$ uniformemente.
  \item Calcula el ``compromiso'': $\vv{w} = \mm{A}\vv{y} \pmod{q}$.
  \item Se genera un desaf\'io $c$ (hash del mensaje y del compromiso).
  \item El firmante responde con $\vv{z} = \vv{y} + c\vv{s}$.
  \item El verificador acepta si $\mm{A}\vv{z} - c\vv{t} = \mm{A}\vv{y}$
    y $\|\vv{z}\|_\infty$ es peque\~no.
\end{enumerate}

El \emph{mecanismo de aborto} descarta $\vv{z}$ si es demasiado grande o no
tiene distribuci\'on uniforme, previniendo fugas del secreto $\vv{s}$.

\subsection{Generaci\'on de claves ML-DSA-65}

\begin{algorithm}[H]
\caption{$\MLDSA.\KeyGen()$}
\label{alg:mldsa_keygen}
\SetKwInOut{Output}{Salida}
\Output{$(\mathtt{sk}, \mathtt{pk})$}
$\xi \xleftarrow{\$} \{0,1\}^{32}$\;
$(\rho, \rho', K) \leftarrow H(\xi)$\;
Generar $\mm{A} \in R_q^{k\times\ell}$ a partir de $\rho$\;
Muestrear $\vv{s}_1 \leftarrow S_\eta^\ell$,\; $\vv{s}_2 \leftarrow S_\eta^k$\tcp*{coeficientes en $[-\eta, \eta]$}
$\vv{t} \leftarrow \mm{A}\vv{s}_1 + \vv{s}_2$\;
$(\vv{t}_1, \vv{t}_0) \leftarrow \mathrm{Power2Round}(\vv{t})$\;
$\mathtt{pk} \leftarrow (\rho, \vv{t}_1)$;\quad
$\mathtt{sk} \leftarrow (\rho, K, \mathrm{tr}, \vv{s}_1, \vv{s}_2, \vv{t}_0)$\;
\Return $(\mathtt{sk}, \mathtt{pk})$
\end{algorithm}

\subsection{Firma}

\begin{algorithm}[H]
\caption{$\MLDSA.\Sign(\mathtt{sk}, m)$}
\label{alg:mldsa_sign}
\SetKwInOut{Output}{Salida}
\Output{Firma $\sigma$}
\Repeat{no abortar}{
  $\vv{y} \leftarrow S_\gamma^{\ell}$\;
  $\vv{w} \leftarrow \mm{A}\vv{y}$\;
  $c \leftarrow H(\mathrm{tr} \| m \| \mathrm{HighBits}(\vv{w}))$\tcp*{desaf\'io hash}
  $\vv{z} \leftarrow \vv{y} + c\vv{s}_1$\;
  \If{$\|\vv{z}\|_\infty \geq \gamma_1 - \beta$
      \textbf{o} $\|\mathrm{LowBits}(\vv{w} - c\vv{s}_2)\|_\infty \geq \gamma_2 - \beta$}{
    \textbf{abortar}\tcp*{reiniciar}
  }
  $\vv{h} \leftarrow \mathrm{MakeHint}(-c\vv{t}_0,\; \vv{w} - c\vv{s}_2 + c\vv{t}_0)$\;
}
\Return $\sigma \leftarrow (\tilde{c}, \vv{z}, \vv{h})$
\end{algorithm}

\subsection{Verificaci\'on}

\begin{algorithm}[H]
\caption{$\MLDSA.\Verify(\mathtt{pk}, m, \sigma)$}
\label{alg:mldsa_verify}
\SetKwInOut{Output}{Salida}
\Output{$\{0,1\}$}
$(\rho, \vv{t}_1) \leftarrow \mathtt{pk}$;\quad
$(\tilde{c}, \vv{z}, \vv{h}) \leftarrow \sigma$\;
\If{$\|\vv{z}\|_\infty \geq \gamma_1 - \beta$}{\Return 0}
$\vv{w}' \leftarrow \mathrm{UseHint}(\vv{h},\; \mm{A}\vv{z} - c\vv{t}_1 \cdot 2^d)$\;
$c' \leftarrow H(\mathrm{tr} \| m \| \vv{w}')$\;
\Return $c' = \tilde{c}$
\end{algorithm}

\subsection{Par\'ametros de ML-DSA-65}

\begin{table}[H]
\centering
\caption{Par\'ametros de ML-DSA-65 (FIPS 204, nivel 3)}
\label{tab:params_mldsa65}
\begin{tabular}{lll}
\toprule
\textbf{Par\'ametro} & \textbf{Valor} & \textbf{Descripci\'on} \\
\midrule
$n$ & 256 & Grado del polinomio c\'iclico \\
$q$ & 8\,380\,417 & M\'odulo primo \\
$(k, \ell)$ & (6, 5) & Dimensi\'on de vectores secreto/respuesta \\
$\eta$ & 4 & Rango del secreto \\
$\gamma_1$ & $2^{19}$ & Rango del vector $\vv{y}$ \\
$\gamma_2$ & $(q-1)/32$ & Radio de redondeo \\
$\beta$ & 120 & Umbral de aborto \\
\midrule
\multicolumn{3}{l}{\textbf{Tama\~nos de estructuras}} \\
Clave p\'ublica & 1\,952 bytes & \\
Clave privada & 4\,032 bytes & \\
Firma & 3\,309 bytes & \\
\bottomrule
\end{tabular}
\end{table}

\section{Comparaci\'on con ECDSA P-256}

\begin{table}[H]
\centering
\caption{ML-DSA-65 vs.\ ECDSA P-256}
\label{tab:comp_firmas}
\begin{tabular}{lll}
\toprule
\textbf{Propiedad} & \textbf{ML-DSA-65} & \textbf{ECDSA P-256} \\
\midrule
Nivel de seguridad & $\sim$192 bits & $\sim$128 bits \\
Base matem\'atica & Module-LWE & Logaritmo discreto EC \\
Resistencia cu\'antica & S\'i (MLWE) & No (algoritmo de Shor) \\
Clave p\'ublica & 1\,952 bytes & $\sim$64 bytes \\
Firma & 3\,309 bytes & $\sim$64 bytes \\
Tiempo de firma & $\sim$0.37 ms & $\sim$0.12 ms \\
Tiempo de verificaci\'on & $\sim$0.21 ms & $\sim$0.18 ms \\
\bottomrule
\end{tabular}
\end{table}

Las firmas ML-DSA son m\'as grandes, pero ofrecen seguridad post-cu\'antica
comprobada con tiempos de operaci\'on comparables.

%###############################################################
\part{El Protocolo de Votaci\'on Electr\'onica}
%###############################################################

%==============================================================
\chapter{Dise\~no Formal del Protocolo}
%==============================================================
\label{chap:protocolo}

\section{Motivaci\'on: por qu\'e PQC en votaci\'on}

Los sistemas de votaci\'on electr\'onica almacenan boletas cifradas que
pueden ser recopiladas hoy y descifradas en el futuro cuando existan
computadoras cu\'anticas potentes (principio ``harvest now, decrypt later'').
Por esta raz\'on, la confidencialidad del voto debe garantizarse
\emph{incluso frente a adversarios cu\'anticos futuros}.

\section{Entidades del protocolo}

\begin{table}[H]
\centering
\caption{Roles del protocolo de votaci\'on}
\label{tab:roles}
\begin{tabular}{clp{7.5cm}}
\toprule
\textbf{ID} & \textbf{Rol} & \textbf{Responsabilidades y claves} \\
\midrule
$V$ & Votante &
  Ciudadano habilitado. Par ML-DSA $(sk_V, pk_V)$ generado en registro.
  Cifra y firma su boleta. \\
\addlinespace
$AR$ & Autoridad de Registro &
  Verifica identidad, emite credenciales firmadas. Par ML-DSA
  $(sk_{AR}, pk_{AR})$. \\
\addlinespace
$SV$ & Servidor de Votaci\'on &
  Recibe boletas, verifica firmas, almacena \emph{sin descifrar}.
  Par ML-KEM $(ek_{SV}, dk_{SV})$ y par ML-DSA $(sk_{SV}, pk_{SV})$. \\
\addlinespace
$AC$ & Autoridad de Conteo &
  Custodia $dk_{SV}$. Desencapsula votos y realiza el escrutinio. \\
\bottomrule
\end{tabular}
\end{table}

\begin{minota}[Separaci\'on de responsabilidades]
  El $SV$ almacena pero \emph{nunca} descifra. La $AC$ descifra pero
  \emph{no} participa en la verificaci\'on de identidad. Esta separaci\'on
  garantiza la confidencialidad del voto durante la jornada electoral.
\end{minota}

\section{Par\'ametros criptogr\'aficos del protocolo}

\begin{table}[H]
\centering
\caption{Primitivas criptogr\'aficas del protocolo}
\label{tab:primitivas}
\begin{tabular}{llll}
\toprule
\textbf{Primitiva} & \textbf{Est\'andar} & \textbf{Nivel seg.} & \textbf{Uso} \\
\midrule
ML-KEM-768 & FIPS 203 & $\sim$AES-192 & Cifrado del voto \\
ML-DSA-65  & FIPS 204 & $\sim$AES-192 & Autenticaci\'on, firma \\
AES-256-GCM & NIST SP 800-38D & 256 bits & Cifrado sim\'etrico \\
SHA3-256 & FIPS 202 & 128 bits colisi\'on & Recibos, auditor\'ia \\
\bottomrule
\end{tabular}
\end{table}

\section{Fases del protocolo}

\subsection{Fase (a) --- Configuraci\'on y generaci\'on de claves}

\begin{enumerate}
  \item $AR$ genera par ML-DSA:
    $(sk_{AR}, pk_{AR}) \leftarrow \MLDSA.\KeyGen()$
  \item $SV$ genera par ML-KEM:
    $(ek_{SV}, dk_{SV}) \leftarrow \MLKEM.\KeyGen()$
  \item $SV$ genera par ML-DSA para recibos:
    $(sk_{SV}, pk_{SV}) \leftarrow \MLDSA.\KeyGen()$
  \item $SV$ entrega $dk_{SV}$ a $AC$ por \emph{canal seguro fuera de banda}.
  \item Se publican $pk_{AR}$, $ek_{SV}$, $pk_{SV}$ en directorio de confianza.
\end{enumerate}

\subsection{Fase (b) --- Registro y autenticaci\'on del Votante}

\begin{enumerate}
  \item $V$ presenta su identidad $ID_V$ ante $AR$.
  \item $AR$ verifica $ID_V$ contra el padr\'on electoral oficial.
  \item Si v\'alido y no registrado:
    \begin{itemize}
      \item $(sk_V, pk_V) \leftarrow \MLDSA.\KeyGen()$
      \item $\mathit{cred}_V = (ID_V \| pk_V \| ts)$
      \item $\sigma_{AR} = \MLDSA.\Sign(sk_{AR}, \mathit{cred}_V)$
      \item Credencial $= (\mathit{cred}_V, \sigma_{AR})$
      \item Registrar $(ID_V, pk_V,\text{``Registrado''})$ en padr\'on.
      \item Entregar $(sk_V, \text{Credencial})$ a $V$ por canal seguro.
    \end{itemize}
\end{enumerate}

\subsection{Fase (c) --- Cifrado del voto y firma de la boleta}

El votante $V$ con $(sk_V, \text{Credencial})$ y conociendo $ek_{SV}$:

\begin{enumerate}
  \item $m = \mathrm{encode}(\text{voto})$
  \item $(K, ct) \leftarrow \MLKEM.\Encaps(ek_{SV})$
  \item $\mathit{voto\_cifrado} = \mathrm{AES\text{-}256\text{-}GCM.Enc}(K, m, \mathit{nonce})$
  \item $\mathit{boleta} = (ct \| \mathit{voto\_cifrado} \| \mathit{nonce}
    \| \text{Credencial} \| ts)$
  \item $\sigma_V = \MLDSA.\Sign(sk_V, \mathit{boleta})$
  \item $\mathit{paquete} = (\mathit{boleta}, \sigma_V)$
\end{enumerate}

\subsection{Fase (d) --- Env\'io al Servidor y verificaci\'on}

\begin{enumerate}
  \item $V \to SV$: $\mathit{paquete}$
  \item $SV$ verifica la credencial: $\MLDSA.\Verify(pk_{AR}, \mathit{cred}_V, \sigma_{AR}) = 1$
  \item $SV$ verifica la firma del votante: $\MLDSA.\Verify(pk_V, \mathit{boleta}, \sigma_V) = 1$
  \item $SV$ verifica que $ID_V$ no ha votado (anti-duplicado).
  \item Si todo es correcto: aceptar. Si falla: rechazar.
\end{enumerate}

\subsection{Fase (e) --- Registro del voto cifrado}

\begin{enumerate}
  \item $SV$ registra $(\mathit{ID}_V, \mathit{boleta})$ en la \'anfora digital.
  \item $SV$ marca $ID_V$ como ``Ya vot\'o''.
  \item $SV$ emite recibo:
    $\mathit{hash} = \mathrm{SHA3\text{-}256}(\mathit{boleta} \| ts_r)$,\quad
    $\sigma_r = \MLDSA.\Sign(sk_{SV}, \mathit{hash})$
  \item $SV \to V$: $(\mathit{hash}, ts_r, \sigma_r)$
\end{enumerate}

\subsection{Fase (f) --- Desencapsulamiento y conteo final}

Con la jornada cerrada, $AC$ obtiene la \'anfora:
\begin{enumerate}
  \item Para cada boleta$_i$:
    \begin{itemize}
      \item $K_i \leftarrow \MLKEM.\Decaps(dk_{SV}, ct_i)$
      \item $m_i = \mathrm{AES\text{-}256\text{-}GCM.Dec}(K_i, \mathit{voto\_cifrado}_i, \mathit{nonce}_i)$
      \item Re-verificar $\sigma_{V_i}$ sobre boleta$_i$.
      \item $\mathit{conteo}[m_i] \mathrel{+}= 1$
    \end{itemize}
  \item $AC$ publica resultado $+$ $\{\mathrm{SHA3\text{-}256}(\mathit{boleta}_i)\}$.
\end{enumerate}

\section{Estructuras de datos}

\subsection{Credencial de votante}

\begin{lstlisting}[language=Python, caption={Credencial de votante}]
Credencial = {
    "cred_V": {
        "ID_V":  str,             # Identidad del votante
        "pk_V":  bytes[1952],     # Clave publica ML-DSA-65
        "ts":    datetime
    },
    "sigma_AR": bytes[3309]       # Firma ML-DSA-65 de la AR
}
\end{lstlisting}

\subsection{Boleta de sufragio}

\begin{lstlisting}[language=Python, caption={Boleta de sufragio}]
Boleta = {
    "ct":           bytes[1088], # ML-KEM-768 ciphertext
    "voto_cifrado": bytes,       # AES-256-GCM(K, m, nonce)
    "nonce":        bytes[12],   # Nonce AES-256-GCM
    "Credencial":   Credencial,
    "ts_boleta":    datetime
}
# Tamano total estimado: ~16,200 bytes (16 KB)
\end{lstlisting}

\subsection{Paquete de voto}

\begin{lstlisting}[language=Python, caption={Paquete de voto final}]
Paquete = {
    "boleta":  Boleta,
    "sigma_V": bytes[3309]  # Firma ML-DSA-65 del votante
}
\end{lstlisting}

%==============================================================
\chapter{Propiedades de Seguridad del Protocolo}
%==============================================================
\label{chap:seguridad}

\section{Modelo de seguridad}

El protocolo fue dise\~nado y verificado bajo el \textbf{modelo de atacante
Dolev-Yao} (1983): el adversario controla completamente la red (intercepta,
modifica, reenv\'ia e inyecta mensajes), pero \emph{no puede romper las
primitivas criptogr\'aficas}.

\section{Propiedad 1: Confidencialidad del voto}

\begin{midef}{Confidencialidad del voto}
  Un adversario polinomial $\mathcal{A}$ que observa todas las boletas no
  puede determinar el contenido del voto $m$ con probabilidad no despreciable:
  \[
    \Pr[\mathcal{A}(\mathit{paquete}) = m] \leq \tfrac{1}{2} + \negl(\lambda).
  \]
\end{midef}

\textbf{Mecanismo:} La seguridad IND-CCA2 de ML-KEM garantiza que $K$ es
computacionalmente indistinguible de aleatorio sin $dk_{SV}$. Como $m$ est\'a
cifrado con AES-256-GCM bajo $K$, tampoco es recuperable. Adem\'as, $dk_{SV}$
es custodiado s\'olo por $AC$.

\section{Propiedad 2: Integridad del voto}

\begin{midef}{Integridad del voto}
  Todo voto contabilizado corresponde exactamente a un voto \emph{emitido}
  por un votante leg\'itimo. Ning\'un voto puede ser modificado sin detecci\'on.
\end{midef}

\textbf{Mecanismo:} La propiedad EUF-CMA de ML-DSA garantiza que ning\'un
adversario puede forjar $\sigma_V$ sin $sk_V$. La $AC$ re-verifica cada firma
antes de contabilizar; si la boleta fue alterada, la verificaci\'on falla.

\section{Propiedad 3: Autenticidad votante-autoridad}

\begin{midef}{Autenticidad}
  Cualquier boleta aceptada por el $SV$ proviene de un votante registrado
  por la $AR$ leg\'itima.
\end{midef}

\textbf{Mecanismo:} La credencial vincula $pk_V$ a $ID_V$ mediante
$\sigma_{AR}$ (EUF-CMA de ML-DSA). El $SV$ verifica ambas firmas.

\section{Propiedad 4: Unicidad del voto}

El $SV$ mantiene estado por $ID_V$. Una vez en estado ``Ya vot\'o'', cualquier
intento posterior de votar con la misma identidad es rechazado.

\section{Propiedad 5: Freshness y anti-replay}

El $ts_{\mathit{boleta}}$ (timestamp) y el $\mathit{nonce}$ (aleatorio) en
cada boleta previenen ataques de repetici\'on: una boleta capturada y
retransmitida tiene el mismo par $(ts, \mathit{nonce})$ que el $SV$ detecta.

\section{Supuestos del modelo}

\begin{enumerate}
  \item \textbf{MLWE hardness:} Module-LWE es dif\'icil para adversarios
    cu\'anticos.
  \item \textbf{Canal fuera de banda:} Entrega de $(sk_V, \text{Credencial})$
    y de $dk_{SV}$ por canales seguros externos.
  \item \textbf{Directorio de confianza:} Claves p\'ublicas distribuidas por
    canal aut\'entico.
  \item \textbf{Dispositivo del votante:} $sk_V$ no est\'a comprometido.
\end{enumerate}

%###############################################################
\part{Implementaci\'on y Evaluaci\'on Emp\'irica}
%###############################################################

%==============================================================
\chapter{Implementaci\'on del Protocolo}
%==============================================================
\label{chap:implementacion}

\section{Arquitectura por capas}

\begin{description}
  \item[Capa C (PQClean)] Operaciones criptogr\'aficas de bajo nivel
    implementadas en C usando las fuentes de referencia del proyecto
    \textbf{PQClean}: portables, auditadas, sin efectos secundarios de
    temporizaci\'on.
  \item[Capa Python] Un wrapper \texttt{ctypes} invoca las funciones C.
    Cada entidad del protocolo tiene su propio m\'odulo con l\'ogica de
    alto nivel.
\end{description}

\section{Motor criptogr\'afico C --- PQClean}

PQClean provee implementaciones de referencia auditadas:
\begin{itemize}
  \item C portable (sin SIMD ni ensamblador espec\'ifico).
  \item Compilable en macOS (ARM64/x86-64) y Linux.
  \item Sin dependencias externas.
  \item Verificado contra vectores KAT del NIST.
\end{itemize}

Biblioteca generada: \texttt{src/pqc/libpqc\_core.dylib} (macOS) o
\texttt{libpqc\_core.so} (Linux).

\begin{lstlisting}[language=C, caption={API de la capa C --- pqc\_core.h}]
/* ML-KEM-768 (FIPS 203) */
int pqc_kem_keypair(uint8_t *ek, uint8_t *dk);
int pqc_kem_enc(uint8_t *ct, uint8_t *ss, const uint8_t *ek);
int pqc_kem_dec(uint8_t *ss, const uint8_t *ct, const uint8_t *dk);

/* ML-DSA-65 (FIPS 204) */
int pqc_sign_keypair(uint8_t *pk, uint8_t *sk);
int pqc_sign(uint8_t *sig, size_t *siglen,
             const uint8_t *msg, size_t msglen,
             const uint8_t *sk);
int pqc_sign_verify(const uint8_t *sig, size_t siglen,
                    const uint8_t *msg, size_t msglen,
                    const uint8_t *pk);
\end{lstlisting}

\section{Wrapper Python (ctypes)}

\begin{lstlisting}[language=Python, caption={Fragmento de pqc\_wrapper.py}]
import ctypes, os

lib = ctypes.CDLL(
    os.path.join(os.path.dirname(__file__), "libpqc_core.dylib")
)

KEM_PK_BYTES = 1184
KEM_SK_BYTES = 2400
KEM_CT_BYTES = 1088
KEM_SS_BYTES = 32

def kem_keypair():
    ek = (ctypes.c_uint8 * KEM_PK_BYTES)()
    dk = (ctypes.c_uint8 * KEM_SK_BYTES)()
    assert lib.pqc_kem_keypair(ek, dk) == 0
    return bytes(ek), bytes(dk)

def kem_encaps(ek: bytes):
    ct = (ctypes.c_uint8 * KEM_CT_BYTES)()
    ss = (ctypes.c_uint8 * KEM_SS_BYTES)()
    ek_buf = (ctypes.c_uint8 * KEM_PK_BYTES).from_buffer_copy(ek)
    assert lib.pqc_kem_enc(ct, ss, ek_buf) == 0
    return bytes(ct), bytes(ss)

def kem_decaps(dk: bytes, ct: bytes):
    ss = (ctypes.c_uint8 * KEM_SS_BYTES)()
    dk_buf = (ctypes.c_uint8 * KEM_SK_BYTES).from_buffer_copy(dk)
    ct_buf = (ctypes.c_uint8 * KEM_CT_BYTES).from_buffer_copy(ct)
    assert lib.pqc_kem_dec(ss, ct_buf, dk_buf) == 0
    return bytes(ss)
\end{lstlisting}

\section{M\'odulos Python del protocolo}

\begin{table}[H]
\centering
\caption{M\'odulos Python del protocolo PQC}
\label{tab:modulos}
\begin{tabular}{lp{9cm}}
\toprule
\textbf{M\'odulo} & \textbf{Responsabilidad} \\
\midrule
\texttt{autoridad\_registro.py} &
  Par ML-DSA de la AR; emisi\'on y firma de credenciales; gesti\'on del padr\'on. \\
\texttt{votante.py} &
  Encapsulamiento ML-KEM, cifrado AES-GCM, construcci\'on y firma de la boleta. \\
\texttt{servidor\_votacion.py} &
  Verificaci\'on doble ML-DSA, anti-duplicado, almacenamiento en \'anfora, recibo. \\
\texttt{autoridad\_conteo.py} &
  Desencapsulamiento ML-KEM, descifrado AES-GCM, re-verificaci\'on, conteo. \\
\texttt{protocolo\_pqc.py} &
  Orquestador de alto nivel: flujo completo de 6 fases. \\
\texttt{crypto\_utils.py} &
  Primitivas sim\'etricas: AES-256-GCM, SHA3-256, KDF (HKDF). \\
\bottomrule
\end{tabular}
\end{table}

\section{Validaci\'on KAT (Known Answer Tests)}

\begin{table}[H]
\centering
\caption{Resultados de validaci\'on KAT --- NIST ACVP}
\label{tab:kat}
\begin{tabular}{llll}
\toprule
\textbf{Algoritmo} & \textbf{Est\'andar} & \textbf{Vectores PASS} & \textbf{Estado} \\
\midrule
ML-KEM-768 & FIPS 203 & 100/100 & \textbf{PASS} \\
ML-DSA-65  & FIPS 204 & 100/100 & \textbf{PASS} \\
\bottomrule
\end{tabular}
\end{table}

%==============================================================
\chapter{Benchmarking y An\'alisis Estad\'istico}
%==============================================================
\label{chap:benchmarks}

\section{Dise\~no experimental}

\begin{table}[H]
\centering
\caption{Dise\~no del experimento}
\label{tab:diseno}
\begin{tabular}{ll}
\toprule
\textbf{Par\'ametro} & \textbf{Valor} \\
\midrule
Escenarios ($N$ votantes) & $\{10,\; 50,\; 100,\; 500,\; 1000\}$ \\
Repeticiones por escenario & 30 \\
Total de corridas & 300 \\
Total de operaciones cripto & $\sim$99\,000 \\
Hardware & macOS 26.5.2, Apple Silicon ARM64 \\
Compilador C & Clang 21.0.0 \\
Python & 3.14.4 \\
\bottomrule
\end{tabular}
\end{table}

\section{Latencia por voto (extremo a extremo)}

\begin{table}[H]
\centering
\caption{Latencia media por voto}
\label{tab:latencia}
\begin{tabular}{lllll}
\toprule
$N$ & Media PQC (ms) & IC 95\% PQC & Media Cl\'asico (ms) & Factor \\
\midrule
10   & 1.333 & [1.281, 1.385] & 0.458 & $2.91\times$ \\
50   & 1.330 & [1.311, 1.348] & 0.457 & $2.91\times$ \\
100  & 1.347 & [1.329, 1.364] & 0.457 & $2.95\times$ \\
500  & 1.351 & [1.333, 1.368] & 0.459 & $2.94\times$ \\
1000 & 1.343 & [1.336, 1.351] & 0.459 & $2.93\times$ \\
\bottomrule
\end{tabular}
\end{table}

\begin{minota}[Interpretaci\'on]
  El protocolo PQC es $\sim$2.9$\times$ m\'as lento que el cl\'asico, pero
  \emph{ambos est\'an muy por debajo del umbral de 200 ms}. El overhead es
  irrelevante para la experiencia del usuario ($<$2\,ms vs.\ 200\,ms).
\end{minota}

\section{Apertura de \'anfora --- Escrutinio}

\begin{table}[H]
\centering
\caption{Desencapsulamiento/descifrado por boleta en escrutinio}
\label{tab:anfora}
\begin{tabular}{lll}
\toprule
\textbf{Operaci\'on} & \textbf{Tiempo} & \textbf{Comparaci\'on} \\
\midrule
ML-KEM-768 Desencapsulamiento & 0.042 ms & --- \\
RSA-OAEP-2048 Descifrado & 92.137 ms & PQC es \textbf{2175$\times$ m\'as r\'apido} \\
\bottomrule
\end{tabular}
\end{table}

Para 10\,000 votos: PQC $\approx$ 420\,ms vs.\ RSA $\approx$ 921\,s (15.4 min).

\section{Micro-benchmarks por operaci\'on at\'omica}

\begin{table}[H]
\centering
\caption{Tiempos por operaci\'on individual (100 repeticiones)}
\label{tab:microbench}
\begin{tabular}{llr}
\toprule
\textbf{Operaci\'on} & \textbf{Algoritmo} & \textbf{Tiempo (ms)} \\
\midrule
KeyGen & ML-KEM-768 & $\sim$0.18 \\
KeyGen & ML-DSA-65 & $\sim$0.21 \\
Encapsulamiento & ML-KEM-768 & $\sim$0.19 \\
Firma & ML-DSA-65 & $\sim$0.37 \\
Verificaci\'on & ML-DSA-65 & $\sim$0.21 \\
\textbf{Desencapsulamiento} & \textbf{ML-KEM-768} & $\sim$\textbf{0.042} \\
\midrule
KeyGen & RSA-2048 & $\sim$28.5 \\
Cifrado & RSA-OAEP-2048 & $\sim$0.18 \\
\textbf{Descifrado} & \textbf{RSA-OAEP-2048} & $\sim$\textbf{92.14} \\
Firma & ECDSA P-256 & $\sim$0.12 \\
Verificaci\'on & ECDSA P-256 & $\sim$0.18 \\
\bottomrule
\end{tabular}
\end{table}

\section{Sobrecarga de comunicaci\'on}

\begin{table}[H]
\centering
\caption{Comparaci\'on de tama\~nos de estructuras}
\label{tab:tamanos}
\begin{tabular}{lll}
\toprule
\textbf{Estructura} & \textbf{PQC} & \textbf{Cl\'asico} \\
\midrule
Clave p\'ublica de elecci\'on & 1\,184 bytes & 294 bytes \\
Firma digital & 3\,309 bytes & 72 bytes \\
Ciphertext por boleta & 1\,088 bytes & 256 bytes \\
\textbf{Boleta total} & \textbf{$\sim$16\,200 bytes} & \textbf{$\sim$1\,200 bytes} \\
\bottomrule
\end{tabular}
\end{table}

La boleta PQC es 13.5$\times$ m\'as grande, pero en red 4G (10 Mbps) s\'olo
a\~nade $\sim$13\,ms de transmisi\'on adicional por boleta --- operativamente
despreciable.

\section{An\'alisis estad\'istico inferencial}

\subsection{Validaci\'on de la hip\'otesis}

\textbf{Hip\'otesis de la tesis:}
\begin{quote}
  \textit{``El protocolo PQC puede operar dentro de los l\'imites de
  viabilidad operativa (latencia $\leq$ 200\,ms por voto) en condiciones
  de carga electoral realistas.''}
\end{quote}

\textbf{Resultado:} $H_0$ rechazada en todos los escenarios con
$p < 0.0001$ (Mann-Whitney U). La latencia media de 1.34\,ms es
\textbf{149 veces menor} al umbral de 200\,ms.

\subsection{Tama\~no del efecto --- Cohen's $d$}

\begin{table}[H]
\centering
\caption{Tama\~no del efecto (Cohen's $d$) por escenario}
\label{tab:cohen}
\begin{tabular}{lll}
\toprule
\textbf{Escenario} & \textbf{Cohen's $d$} & \textbf{Interpretaci\'on} \\
\midrule
$N=10$   & 7.23  & Efecto enorme ($>$0.8) \\
$N=50$   & 18.45 & Efecto enorme \\
$N=100$  & 19.31 & Efecto enorme \\
$N=500$  & 18.92 & Efecto enorme \\
$N=1000$ & 26.71 & Efecto enorme \\
\bottomrule
\end{tabular}
\end{table}

Valores de Cohen's $d > 7$ confirman que la diferencia es
\emph{sustancial y no un artefacto estad\'istico}.

%###############################################################
\part{Verificaci\'on Formal y Conclusiones}
%###############################################################

%==============================================================
\chapter{Verificaci\'on Formal con Tamarin Prover}
%==============================================================
\label{chap:tamarin}

\begin{minota}[Sobre este cap\'itulo]
  La gu\'ia completa de Tamarin Prover existe como documento separado
  (\texttt{tamarin/guia\_tamarin\_libro.tex}). Este cap\'itulo presenta
  s\'olo el contexto y los resultados de la verificaci\'on formal del
  protocolo de votaci\'on.
\end{minota}

\section{\textquestiondown Para qu\'e sirve la verificaci\'on formal?}

La verificaci\'on formal permite demostrar matem\'aticamente (no s\'olo
testear) que un protocolo cumple sus propiedades de seguridad bajo
\emph{todas las posibles ejecuciones}, incluyendo las del atacante.

Tamarin Prover es una herramienta de verificaci\'on autom\'atica de protocolos
de seguridad basada en el c\'alculo de reescritura de multiconjuntos. Trabaja
bajo el \textbf{modelo Dolev-Yao}: adversario controla la red pero no puede
romper las primitivas criptogr\'aficas.

\section{Abstracci\'on criptogr\'afica en el modelo}

ML-KEM y ML-DSA se modelan como cajas negras perfectas (funciones
simb\'olicas), sin modelar la matem\'atica de ret\'iculos:

\begin{itemize}
  \item \textbf{ML-KEM:}
    $\mathtt{decaps}(dk,\; \mathtt{encaps}(ek, r)) = r$
  \item \textbf{ML-DSA:}
    $\mathtt{verify}(\mathtt{sign}(m, sk),\; m,\; \mathtt{pk}(sk)) = \mathtt{true}$
  \item \textbf{AEAD:}
    $\mathtt{sdec}(\mathtt{senc}(m, k),\; k) = m$
\end{itemize}

\section{Lemas verificados}

\begin{table}[H]
\centering
\caption{Resultados de verificaci\'on formal con Tamarin Prover}
\label{tab:tamarin}
\begin{tabular}{llll}
\toprule
\textbf{Lema} & \textbf{Tipo} & \textbf{Propiedad} & \textbf{Estado} \\
\midrule
\texttt{executability} & \texttt{exists-trace} &
  Sanity check: traza v\'alida completa & \textbf{VERIFIED} \\
\texttt{secreto\_voto} & \texttt{all-traces} &
  Adversario no aprende $m$ & \textbf{VERIFIED} \\
\texttt{integridad\_voto} & \texttt{all-traces} &
  Correspondencia emitido/contabilizado & \textbf{VERIFIED} \\
\texttt{autenticidad\_votante} & \texttt{all-traces} &
  Injective agreement, previene replay & \textbf{VERIFIED} \\
\bottomrule
\end{tabular}
\end{table}

Los cuatro lemas se verificaron en menos de 5 segundos cada uno.

%==============================================================
\chapter{Suite de Pruebas y Calidad del C\'odigo}
%==============================================================
\label{chap:pruebas}

\section{Suite de pruebas automatizadas (pytest)}

\begin{table}[H]
\centering
\caption{Suite de pruebas automatizadas}
\label{tab:tests}
\begin{tabular}{lll}
\toprule
\textbf{Archivo} & \textbf{Cobertura} & \textbf{Resultado} \\
\midrule
\texttt{test\_pqc.py} & Protocolo PQC completo & 8/8 PASS \\
\texttt{test\_clasico.py} & Protocolo cl\'asico & 7/7 PASS \\
\texttt{test\_web\_api.py} & Endpoints Flask & 5/5 PASS \\
\midrule
\textbf{Total} & & \textbf{20/20 PASS} \\
\bottomrule
\end{tabular}
\end{table}

\section{M\'etricas de calidad}

\begin{table}[H]
\centering
\caption{M\'etricas de calidad del artefacto}
\label{tab:calidad}
\begin{tabular}{ll}
\toprule
\textbf{M\'etrica} & \textbf{Valor} \\
\midrule
Tests unitarios/integraci\'on & 20/20 (100\% PASS) \\
Vectores KAT NIST validados & 200/200 (100\% PASS) \\
Lemas formales verificados & 4/4 (100\% VERIFIED) \\
L\'ineas de c\'odigo Python & $\sim$2\,500 \\
L\'ineas de c\'odigo C & $\sim$500 \\
Commits en el repositorio & 8+ commits \\
\bottomrule
\end{tabular}
\end{table}

%==============================================================
\chapter{Conclusiones y Trabajo Futuro}
%==============================================================
\label{chap:conclusiones}

\section{S\'intesis de resultados}

Este proyecto dise\~n\'o, implement\'o y evalu\'o emp\'iricamente un protocolo
de votaci\'on electr\'onica resistente a amenazas cu\'anticas empleando los
est\'andares ML-KEM-768 (FIPS 203) y ML-DSA-65 (FIPS 204) del NIST.

\subsection{Respecto a la hip\'otesis}

La hip\'otesis de investigaci\'on est\'a \textbf{formalmente confirmada} con
$\alpha = 0.0001$:
\begin{itemize}
  \item Latencia media: \textbf{1.34\,ms por voto} $\ll$ umbral de 200\,ms.
  \item Diferencia estad\'isticamente significativa ($p < 0.0001$, Mann-Whitney U).
  \item Efecto enorme (Cohen's $d > 7$ en todos los escenarios).
\end{itemize}

\subsection{Ventajas del protocolo PQC}

\begin{enumerate}
  \item \textbf{Resistencia cu\'antica comprobada:} ML-KEM-768 y ML-DSA-65
    basan su seguridad en Module-LWE, intractable para computadoras cu\'anticas.
  \item \textbf{Escrutinio ultrarrápido:} Desencapsulamiento ML-KEM-768 es
    2\,175$\times$ m\'as r\'apido que RSA-OAEP en apertura de \'anfora.
  \item \textbf{Estandarizaci\'on oficial:} NIST FIPS 203/204 garantizan
    interoperabilidad y soporte a largo plazo.
  \item \textbf{Seguridad formal verificada:} 4 lemas demostrados con
    Tamarin Prover.
  \item \textbf{Implementaci\'on auditada:} Basada en PQClean; validada con
    200 vectores KAT oficiales del NIST.
\end{enumerate}

\section{Limitaciones}

\begin{enumerate}
  \item \textbf{Sobrecarga de comunicaci\'on:} Boletas PQC son 13.5$\times$
    m\'as grandes. En redes $<$1\,Mbps podr\'ia ser relevante.
  \item \textbf{Escalabilidad:} Pruebas hasta $N=1\,000$ votantes. Para
    elecciones nacionales se requiere arquitectura distribuida.
  \item \textbf{Privacidad avanzada:} No implementa mix-nets ni pruebas de
    conocimiento cero para anonimato matem\'aticamente m\'as fuerte.
\end{enumerate}

\section{Trabajo futuro}

\begin{itemize}
  \item Evaluar \textbf{FALCON-512} (firmas compactas, $\sim$666 bytes) como
    alternativa a ML-DSA-65.
  \item Implementar \textbf{esquemas de umbral} para distribuir la custodia de
    $dk_{SV}$ entre m\'ultiples autoridades.
  \item Integrar \textbf{pruebas de conocimiento cero} (ZKP) para verificaci\'on
    de votos sin revelaci\'on del contenido.
  \item Escalar a arquitectura distribuida y evaluar con millones de votantes.
  \item Migrar credenciales a \textbf{FIPS 205 (SLH-DSA)} para larga duraci\'on.
\end{itemize}

%=============================================================
\appendix
%=============================================================

%==============================================================
\chapter{Notaci\'on Matem\'atica}
%==============================================================
\label{app:notacion}

\begin{table}[H]
\centering
\caption{Notaci\'on utilizada en el libro}
\label{tab:notacion}
\begin{tabular}{lp{9cm}}
\toprule
\textbf{S\'imbolo} & \textbf{Significado} \\
\midrule
$\Z$ & Conjunto de los enteros \\
$\Z_q$ & Enteros m\'odulo $q$, $\{0,1,\ldots,q-1\}$ \\
$\Z_q[x]/(x^n+1)$ & Anillo de polinomios c\'iclicos \\
$R_q$ & $\Z_q[x]/(x^n+1)$, abreviatura \\
$R_q^k$ & M\'odulo de rango $k$ sobre $R_q$ \\
$\vv{v}$ & Vector (negrita min\'uscula) \\
$\mm{M}$ & Matriz (negrita may\'uscula) \\
$\|\vv{v}\|$ & Norma euclidiana \\
$\|\vv{v}\|_\infty$ & Norma infinito: $\max_i |v_i|$ \\
$\langle\vv{a},\vv{b}\rangle$ & Producto interior \\
$a \xleftarrow{\$} S$ & $a$ muestreado uniformemente de $S$ \\
$a \leftarrow \chi$ & $a$ muestreado de la distribuci\'on $\chi$ \\
$\negl(\lambda)$ & Funci\'on despreciable en $\lambda$ \\
$\lambda$ & Par\'ametro de seguridad \\
$\lfloor x \rceil$ & Redondeo al entero m\'as cercano \\
$H(\cdot)$ & Funci\'on hash (SHA3/SHAKE) \\
$\|$ & Concatenaci\'on de bytes \\
\bottomrule
\end{tabular}
\end{table}

%==============================================================
\chapter{Glosario de T\'erminos}
%==============================================================
\label{app:glosario}

\begin{description}
  \item[Algoritmo de Shor] Algoritmo cu\'antico (1994) que resuelve la
    factorizaci\'on de enteros y el logaritmo discreto en tiempo polinomial.
    Amenaza a RSA, DH y ECDSA.

  \item[\'Anfora digital] Base de datos de boletas cifradas gestionada por
    el Servidor de Votaci\'on. No puede ser descifrada sin $dk_{SV}$.

  \item[Boleta] Estructura de datos firmada y cifrada que contiene el voto,
    la credencial y metadatos de integridad.

  \item[Credencial] Certificado emitido por la AR que vincula $ID_V$ con
    $pk_V$, firmado con ML-DSA.

  \item[CVP (Closest Vector Problem)] Encontrar el vector m\'as cercano a
    un punto objetivo en un ret\'iculo. NP-dif\'icil.

  \item[EUF-CMA] Existential Unforgeability under Chosen-Message Attack.
    Propiedad de seguridad m\'axima para firmas digitales.

  \item[IND-CCA2] Indistinguishability under Adaptive Chosen-Ciphertext Attack.
    Propiedad de seguridad m\'axima para cifrado p\'ublico. ML-KEM la satisface.

  \item[KEM (Key Encapsulation Mechanism)] Protocolo que permite establecer
    un secreto compartido sin transmitirlo directamente.

  \item[LWE (Learning With Errors)] Problema: dado un sistema lineal con error
    peque\~no, recuperar el secreto o distinguirlo de uniforme. NP-dif\'icil.

  \item[MLWE (Module-LWE)] Variante de LWE sobre m\'odulos de anillos de
    polinomios. Base de ML-KEM y ML-DSA.

  \item[NTT (Number Theoretic Transform)] Transformada an\'aloga a la DFT en
    aritm\'etica modular. Permite multiplicar polinomios en $O(n\log n)$.

  \item[PQC (Post-Quantum Cryptography)] Algoritmos cuya seguridad se
    cree intractable incluso para computadoras cu\'anticas.

  \item[PQClean] Proyecto open-source con implementaciones C de referencia
    de los algoritmos NIST PQC.

  \item[Replay attack] Ataque donde el adversario retransmite un mensaje
    v\'alido capturado para efectos no deseados.

  \item[RLWE (Ring-LWE)] Variante de LWE sobre anillos c\'iclicos $\Z_q[x]/(x^n+1)$.
    M\'as eficiente que LWE est\'andar.

  \item[SIS (Short Integer Solution)] Dado $\mm{A}$, encontrar $\vv{z}$
    corto tal que $\mm{A}\vv{z} = \vv{0} \pmod{q}$. Base de ML-DSA.

  \item[SVP (Shortest Vector Problem)] Encontrar el vector m\'as corto no
    nulo en un ret\'iculo. NP-dif\'icil.

  \item[Transformaci\'on Fujisaki-Okamoto (FO)] T\'ecnica que convierte un
    esquema OW-CPA en IND-CCA2. Usada internamente en ML-KEM.
\end{description}

%==============================================================
\chapter{Flujo del Protocolo en Notaci\'on Alice-Bob}
%==============================================================
\label{app:alice_bob}

\begin{lstlisting}[caption={Protocolo de votaci\'on en notaci\'on Alice-Bob}]
-- FASE (a): Configuracion --
AR  : (sk_AR, pk_AR) <- ML-DSA.KeyGen()
SV  : (ek_SV, dk_SV) <- ML-KEM.KeyGen()
SV  : (sk_SV, pk_SV) <- ML-DSA.KeyGen()
SV  --(seguro)--> AC : dk_SV
AR, SV --> Directorio : {pk_AR, ek_SV, pk_SV}

-- FASE (b): Registro --
V   --> AR : ID_V
AR  : (sk_V, pk_V) <- ML-DSA.KeyGen()
AR  : cred_V = (ID_V || pk_V || ts)
AR  : sigma_AR = ML-DSA.Sign(sk_AR, cred_V)
AR  --(seguro)--> V : (sk_V, (cred_V, sigma_AR))

-- FASE (c): Cifrado del voto --
V   : m = encode(voto)
V   : (K, ct) <- ML-KEM.Encaps(ek_SV)
V   : voto_cifrado = AES-GCM.Enc(K, m, nonce)
V   : boleta = ct || voto_cifrado || nonce
             || (cred_V, sigma_AR) || ts_boleta
V   : sigma_V = ML-DSA.Sign(sk_V, boleta)
V   : paquete = (boleta, sigma_V)

-- FASES (d)-(e): Envio y registro --
V   --> SV : paquete
SV  : assert ML-DSA.Verify(pk_AR, cred_V, sigma_AR)
SV  : assert ML-DSA.Verify(pk_V, boleta, sigma_V)
SV  : assert padron[ID_V].estado != "Ya voto"
SV  : anfora[ID_V] = boleta
SV  : padron[ID_V].estado = "Ya voto"
SV  --> V  : (SHA3-256(boleta||ts_r), sigma_SV)

-- FASE (f): Escrutinio --
SV  --(seguro)--> AC : {anfora}
para cada boleta_i en anfora:
  AC : K_i <- ML-KEM.Decaps(dk_SV, ct_i)
  AC : m_i = AES-GCM.Dec(K_i, voto_cifrado_i)
  AC : assert ML-DSA.Verify(pk_V_i, boleta_i, sigma_V_i)
  AC : conteo[m_i] += 1
AC  --> Publico : (conteo, {SHA3-256(boleta_i)})
\end{lstlisting}

%==============================================================
\chapter{Comparaci\'on PQC vs.\ Cl\'asico}
%==============================================================
\label{app:comparacion}

\begin{table}[H]
\centering
\caption{Comparaci\'on completa: protocolo PQC vs.\ Cl\'asico}
\label{tab:comp_completa}
\begin{tabular}{lll}
\toprule
\textbf{Aspecto} & \textbf{PQC (ML-KEM + ML-DSA)} & \textbf{Cl\'asico (RSA + ECDSA)} \\
\midrule
\multicolumn{3}{l}{\textit{Seguridad}} \\
Resistencia cu\'antica & S\'i (MLWE) & No (algoritmo de Shor) \\
Base matem\'atica & Module-LWE & Factorizaci\'on / DLP \\
Nivel de seguridad & $\sim$192 bits (Cat.\ 3) & $\sim$128 bits \\
Estandarizaci\'on & NIST FIPS 203/204 (2024) & RSA/ECDSA est\'andar \\
\midrule
\multicolumn{3}{l}{\textit{Rendimiento}} \\
Latencia por voto & 1.34\,ms & 0.46\,ms \\
Apertura de \'anfora & 0.042\,ms & 92.14\,ms \\
KeyGen KEM/cifrado & $\sim$0.2\,ms & $\sim$28.5\,ms (RSA) \\
\midrule
\multicolumn{3}{l}{\textit{Tama\~nos}} \\
Clave p\'ublica & 1\,184 bytes (KEM) & 294 bytes (RSA) \\
Firma & 3\,309 bytes & 72 bytes (ECDSA) \\
Boleta total & $\sim$16\,200 bytes & $\sim$1\,200 bytes \\
\midrule
\multicolumn{3}{l}{\textit{Verificaci\'on formal}} \\
Tamarin (lemas) & 4/4 VERIFIED & N/A \\
KAT NIST & 200/200 PASS & N/A \\
\bottomrule
\end{tabular}
\end{table}

%=============================================================
\backmatter
%=============================================================

\chapter*{Referencias}
\addcontentsline{toc}{chapter}{Referencias}

\begin{enumerate}
  \item NIST. \textit{FIPS 203: Module-Lattice-Based Key-Encapsulation
    Mechanism Standard}. 2024.
    \url{https://doi.org/10.6028/NIST.FIPS.203}

  \item NIST. \textit{FIPS 204: Module-Lattice-Based Digital Signature
    Standard}. 2024.
    \url{https://doi.org/10.6028/NIST.FIPS.204}

  \item Regev, O. \textit{On Lattices, Learning with Errors, Random Linear
    Codes, and Cryptography}. JACM, 56(6), 2009.

  \item Lyubashevsky, V., Peikert, C., Regev, O. \textit{On Ideal Lattices
    and Learning with Errors over Rings}. JACM, 60(6), 2013.

  \item Langlois, A., Stehl\'e, D. \textit{Worst-Case to Average-Case
    Reductions for Module Lattices}. DCC, 75(3), 2015.

  \item Ducas, L., Kiltz, E., Lepoint, T., Lyubashevsky, V., Schwabe, P.,
    Seiler, G., Stehl\'e, D. \textit{CRYSTALS-Dilithium: A Lattice-Based
    Digital Signature Scheme}. IACR TCHES, 2018(1), 2018.

  \item Avanzi, R. et al. \textit{CRYSTALS-Kyber Algorithm Specifications}.
    NIST PQC Submission Round 3, 2021.

  \item Fujisaki, E., Okamoto, T. \textit{Secure Integration of Asymmetric
    and Symmetric Encryption Schemes}. JoC, 26(1), 2013.

  \item Shor, P.W. \textit{Polynomial-Time Algorithms for Prime Factorization
    and Discrete Logarithms on a Quantum Computer}. SIAM JoC, 26(5), 1997.

  \item Diffie, W., Hellman, M.E. \textit{New Directions in Cryptography}.
    IEEE TIT, 22(6), 1976.

  \item Lenstra, A.K., Lenstra, H.W., Lov\'asz, L. \textit{Factoring
    Polynomials with Rational Coefficients}. MA, 261(4), 1982. (Algoritmo LLL)

  \item Schmidt, B., Meier, S., Cremers, C., Basin, D. \textit{Automated
    Analysis of Diffie-Hellman Protocols and Advanced Security Properties}.
    IEEE CSF, 2012.

  \item Dolev, D., Yao, A. \textit{On the Security of Public Key Protocols}.
    IEEE TIT, 29(2), 1983.

  \item PQClean Project. \textit{Clean, portable implementations of PQC}.
    GitHub, 2024. \url{https://github.com/PQClean/PQClean}

  \item Vaishnavi, V.K., Kuechler, W. \textit{Design Science Research Methods
    and Patterns}. CRC Press, 2004.

  \item Bernstein, D.J., Lange, T. \textit{Post-Quantum Cryptography}.
    Nature, 549(7671), 2017.
\end{enumerate}

\printindex

\end{document}
